\documentclass[12pt,a4paper]{article}

% ==================== GÓI CƠ BẢN ====================
\usepackage[utf8]{inputenc}
\usepackage[T5]{fontenc}
\usepackage[vietnamese]{babel}
\usepackage{geometry}
\usepackage{parskip}
\usepackage{enumitem}
\usepackage{longtable}
\usepackage{array}
\usepackage{xcolor}
\usepackage{hyperref}
\usepackage{fancyhdr}
\usepackage{lastpage}

% ==================== ĐỊNH DẠNG TRANG ====================
\geometry{
    top=2cm,
    bottom=2cm,
    left=2.5cm,
    right=2cm
}

\setlength{\parindent}{0pt}
\setlength{\parskip}{0.7em}
\setlist[itemize]{leftmargin=1.8em, itemsep=0.35em}
\setlist[enumerate]{leftmargin=1.8em, itemsep=0.35em}

\definecolor{ailmsblue}{RGB}{25, 84, 166}
\definecolor{ailmsgray}{RGB}{90, 90, 90}

\hypersetup{
    colorlinks=true,
    linkcolor=ailmsblue,
    urlcolor=ailmsblue,
    citecolor=ailmsblue,
    pdftitle={Chính sách AI, Dữ liệu và Bảo mật AILMS},
    pdfauthor={Phòng Công nghệ và Dữ liệu AILMS}
}

% ==================== HEADER / FOOTER ====================
\pagestyle{fancy}
\fancyhf{}
\fancyhead[L]{\textcolor{ailmsgray}{AILMS}}
\fancyhead[R]{\textcolor{ailmsgray}{Chính sách AI, Dữ liệu \& Bảo mật}}
\fancyfoot[C]{Trang \thepage/\pageref{LastPage}}
\renewcommand{\headrulewidth}{0.4pt}
\renewcommand{\footrulewidth}{0.4pt}

% ==================== THÔNG TIN TÀI LIỆU ====================
\title{
    \textbf{\LARGE CHÍNH SÁCH SỬ DỤNG TRÍ TUỆ NHÂN TẠO,\\
    QUẢN TRỊ DỮ LIỆU VÀ BẢO MẬT THÔNG TIN AILMS}
}
\author{Phòng Công nghệ và Dữ liệu AILMS}
\date{Phiên bản 1.0 -- Ngày hiệu lực: \today}

\begin{document}

\maketitle

\begin{center}
    \textit{Tài liệu này quy định các nguyên tắc sử dụng trí tuệ nhân tạo, xử lý dữ liệu,
    bảo vệ thông tin và kiểm soát truy cập trên hệ thống AI Personalized Learning
    Management System (AILMS).}
\end{center}

\vspace{1cm}

\begin{longtable}{|>{\bfseries}p{4cm}|p{9.5cm}|}
    \hline
    Tên tài liệu &
    Chính sách sử dụng Trí tuệ nhân tạo, Quản trị dữ liệu và Bảo mật thông tin AILMS \\
    \hline
    Phiên bản & 1.0 \\
    \hline
    Đơn vị quản lý & Phòng Công nghệ và Dữ liệu AILMS \\
    \hline
    Đối tượng áp dụng &
    Học viên, giảng viên, trợ giảng, nhân sự, quản trị viên,
    đối tác và các bên được cấp quyền truy cập hệ thống \\
    \hline
    Chu kỳ rà soát & Tối thiểu 12 tháng một lần hoặc khi có thay đổi quan trọng \\
    \hline
\end{longtable}

\tableofcontents
\newpage

% =========================================================
\section{Mục đích}

Chính sách này được ban hành nhằm thiết lập khuôn khổ thống nhất cho việc phát triển,
triển khai và sử dụng các chức năng trí tuệ nhân tạo, dữ liệu học tập và các thành phần
bảo mật trong hệ thống AILMS.

Chính sách hướng đến các mục tiêu sau:

\begin{itemize}
    \item Bảo đảm AI được sử dụng đúng mục đích, có trách nhiệm và minh bạch.
    \item Bảo vệ dữ liệu cá nhân, dữ liệu học tập và dữ liệu nội bộ của tổ chức.
    \item Ngăn chặn truy cập trái phép, rò rỉ dữ liệu và lạm dụng chức năng AI.
    \item Xác định rõ trách nhiệm của người dùng và đơn vị vận hành hệ thống.
    \item Duy trì khả năng kiểm tra, truy vết và xử lý sự cố.
    \item Tuân thủ các quy định pháp luật hiện hành về bảo vệ dữ liệu và an toàn thông tin.
\end{itemize}

% =========================================================
\section{Phạm vi áp dụng}

Chính sách này áp dụng đối với:

\begin{itemize}
    \item Website, ứng dụng và API thuộc hệ thống AILMS.
    \item Dịch vụ Backend, AI Service, cơ sở dữ liệu, bộ nhớ đệm và kho lưu trữ tệp.
    \item Các chức năng AI Tutor, chatbot tư vấn, gợi ý khóa học, tạo nội dung và RAG.
    \item Dữ liệu do học viên, giảng viên, nhân sự hoặc quản trị viên cung cấp.
    \item Các bên thứ ba được phép tích hợp hoặc xử lý dữ liệu thay mặt AILMS.
\end{itemize}

Việc truy cập hoặc sử dụng hệ thống đồng nghĩa với việc người dùng đồng ý tuân thủ
các quy định được nêu trong chính sách này và các điều khoản liên quan khác của AILMS.

% =========================================================
\section{Giải thích thuật ngữ}

Trong phạm vi tài liệu này, các thuật ngữ dưới đây được hiểu như sau:

\begin{itemize}
    \item \textbf{AILMS:} Hệ thống quản lý học tập có tích hợp trí tuệ nhân tạo và
    cá nhân hóa quá trình học tập.

    \item \textbf{AI Tutor:} Trợ lý học tập sử dụng AI để giải thích kiến thức,
    gợi ý nội dung, hỗ trợ ôn tập và định hướng lộ trình học.

    \item \textbf{AI tạo sinh:} Mô hình có khả năng tạo văn bản, câu hỏi, bài tập,
    bản tóm tắt hoặc nội dung mới dựa trên dữ liệu đầu vào.

    \item \textbf{RAG:} Cơ chế truy xuất dữ liệu liên quan từ kho tri thức trước khi
    mô hình AI tạo câu trả lời.

    \item \textbf{Dữ liệu cá nhân:} Thông tin gắn với hoặc có khả năng xác định một
    cá nhân cụ thể.

    \item \textbf{Dữ liệu học tập:} Thông tin về khóa học, tiến độ, kết quả,
    bài nộp, thời gian học và hành vi tương tác của học viên.

    \item \textbf{RBAC:} Mô hình kiểm soát truy cập dựa trên vai trò của người dùng.

    \item \textbf{Audit Log:} Nhật ký ghi nhận các thao tác quan trọng phục vụ
    kiểm tra, truy vết và điều tra sự cố.
\end{itemize}

% =========================================================
\section{Nguyên tắc sử dụng trí tuệ nhân tạo}

AILMS triển khai AI theo nhiều cấp độ phù hợp với từng trường hợp sử dụng:

\begin{enumerate}
    \item \textbf{Cấp độ 1 -- Rule-based:}
    xử lý câu hỏi hoặc yêu cầu theo các quy tắc và kịch bản được xác định trước.

    \item \textbf{Cấp độ 2 -- Tìm kiếm và xếp hạng:}
    hỗ trợ tìm kiếm khóa học, nội dung và tài nguyên dựa trên từ khóa,
    mức độ liên quan và bộ lọc.

    \item \textbf{Cấp độ 3 -- Machine Learning Recommendation:}
    gợi ý khóa học, lộ trình hoặc nội dung dựa trên hồ sơ và hoạt động học tập.

    \item \textbf{Cấp độ 4 -- AI tạo sinh kết hợp RAG:}
    tạo câu trả lời hoặc nội dung dựa trên mô hình ngôn ngữ và dữ liệu được phép truy xuất.
\end{enumerate}

Việc sử dụng AI phải tuân theo các nguyên tắc:

\begin{itemize}
    \item Có mục đích rõ ràng và phục vụ trực tiếp hoạt động học tập hoặc vận hành.
    \item Không sử dụng AI để đưa ra quyết định có ảnh hưởng nghiêm trọng tới người dùng
    mà không có cơ chế kiểm tra hoặc xác nhận của con người.
    \item Không tự động coi nội dung do AI tạo ra là nguồn thông tin tuyệt đối chính xác.
    \item Không cho phép AI vượt qua cơ chế phân quyền hoặc truy cập dữ liệu trái phép.
    \item Hạn chế tối đa việc đưa dữ liệu cá nhân không cần thiết vào câu lệnh gửi tới AI.
\end{itemize}

% =========================================================
\section{Vai trò và giới hạn của trợ lý AI}

AI Tutor đóng vai trò hỗ trợ, không thay thế hoàn toàn giảng viên, cố vấn,
nhân sự hoặc cán bộ có thẩm quyền.

AI có thể được sử dụng để:

\begin{itemize}
    \item Giải thích khái niệm và nội dung bài học.
    \item Tóm tắt tài liệu hoặc gợi ý tài nguyên học tập.
    \item Gợi ý khóa học và lộ trình dựa trên nhu cầu của học viên.
    \item Tạo bản nháp câu hỏi, bài tập hoặc nội dung giảng dạy.
    \item Hỗ trợ tìm kiếm và truy xuất thông tin trong phạm vi được cấp quyền.
\end{itemize}

AI không được sử dụng làm căn cứ duy nhất để:

\begin{itemize}
    \item Đưa ra quyết định kỷ luật, đình chỉ hoặc khóa tài khoản.
    \item Chấm điểm cuối cùng đối với bài làm có tính chủ quan cao.
    \item Xác định gian lận học thuật mà không có bước kiểm tra của con người.
    \item Đưa ra cam kết pháp lý, tài chính hoặc tuyển dụng thay mặt tổ chức.
    \item Cung cấp lời khuyên chuyên môn trong lĩnh vực y tế, pháp lý hoặc tài chính.
\end{itemize}

% =========================================================
\section{Minh bạch đối với nội dung do AI tạo ra}

AILMS thực hiện các biện pháp minh bạch sau:

\begin{itemize}
    \item Thông báo rõ khi người dùng đang tương tác với AI.
    \item Gắn nhãn phù hợp đối với nội dung được tạo hoặc đề xuất bởi AI.
    \item Cảnh báo rằng câu trả lời của AI có thể chưa chính xác hoặc chưa đầy đủ.
    \item Cho phép người dùng phản hồi, báo cáo hoặc yêu cầu kiểm tra nội dung AI.
    \item Lưu lại thông tin kỹ thuật cần thiết để truy vết quá trình tạo câu trả lời,
    trong phạm vi phù hợp với chính sách bảo mật.
\end{itemize}

Người dùng có trách nhiệm kiểm tra lại các thông tin quan trọng trước khi sử dụng
cho hoạt động học tập, đánh giá hoặc ra quyết định.

% =========================================================
\section{Giới hạn sử dụng và chống lạm dụng AI}

Để bảo đảm hiệu năng, chi phí và an toàn hệ thống, AILMS có thể áp dụng:

\begin{itemize}
    \item Giới hạn số lượng yêu cầu trong một khoảng thời gian.
    \item Giới hạn độ dài câu hỏi, tệp tải lên và lịch sử hội thoại.
    \item Giới hạn số lượt tạo lại câu trả lời.
    \item Giới hạn theo tài khoản, vai trò, địa chỉ IP hoặc gói dịch vụ.
    \item Tạm dừng chức năng AI khi phát hiện hành vi tự động hóa hoặc lạm dụng.
\end{itemize}

Người dùng không được:

\begin{itemize}
    \item Cố gắng vượt qua giới hạn sử dụng hoặc cơ chế phân quyền.
    \item Thực hiện prompt injection nhằm lấy chỉ dẫn hệ thống hoặc dữ liệu nội bộ.
    \item Yêu cầu AI tiết lộ thông tin của người dùng khác.
    \item Dùng hệ thống để tạo mã độc, nội dung lừa đảo hoặc nội dung vi phạm pháp luật.
    \item Gửi hàng loạt yêu cầu tự động gây ảnh hưởng đến tính ổn định của hệ thống.
\end{itemize}

% =========================================================
\section{Thu thập và xử lý dữ liệu}

AILMS có thể thu thập các nhóm dữ liệu sau:

\subsection{Dữ liệu tài khoản}

Bao gồm họ tên, tên đăng nhập, địa chỉ thư điện tử, số điện thoại,
vai trò, trạng thái tài khoản và thông tin xác thực cần thiết.

\subsection{Dữ liệu học tập}

Bao gồm:

\begin{itemize}
    \item Khóa học và gói đào tạo đã đăng ký.
    \item Tiến độ hoàn thành bài học.
    \item Điểm số, bài nộp, bài kiểm tra và phản hồi của giảng viên.
    \item Thời gian truy cập, thời lượng học và lịch sử tương tác.
    \item Lớp học, phiên tư vấn và hoạt động thảo luận.
\end{itemize}

\subsection{Dữ liệu hội thoại AI}

Bao gồm câu hỏi, câu trả lời, phản hồi đánh giá, tệp đính kèm,
thông tin ngữ cảnh và mã định danh cuộc hội thoại.

\subsection{Dữ liệu kỹ thuật}

Bao gồm địa chỉ IP, loại thiết bị, trình duyệt, thời điểm truy cập,
nhật ký hệ thống và thông tin lỗi cần thiết cho việc vận hành và bảo mật.

% =========================================================
\section{Mục đích xử lý dữ liệu}

Dữ liệu được xử lý cho các mục đích:

\begin{itemize}
    \item Cung cấp và duy trì chức năng của hệ thống.
    \item Xác thực tài khoản và kiểm soát truy cập.
    \item Theo dõi tiến độ và quản lý hoạt động học tập.
    \item Cá nhân hóa nội dung, khóa học và lộ trình.
    \item Cải thiện chất lượng tìm kiếm và trợ lý AI.
    \item Phát hiện gian lận, lạm dụng hoặc nguy cơ bảo mật.
    \item Phân tích hiệu năng và khắc phục lỗi kỹ thuật.
    \item Thực hiện nghĩa vụ theo quy định pháp luật hoặc yêu cầu hợp lệ của cơ quan có thẩm quyền.
\end{itemize}

AILMS không sử dụng dữ liệu cá nhân cho mục đích không liên quan nếu chưa có
căn cứ phù hợp hoặc sự đồng ý cần thiết của người dùng.

% =========================================================
\section{Cá nhân hóa hoạt động học tập}

AILMS có thể sử dụng dữ liệu học tập để:

\begin{itemize}
    \item Đề xuất khóa học hoặc bài học tiếp theo.
    \item Phát hiện nội dung người học chưa nắm vững.
    \item Điều chỉnh mức độ khó của tài liệu hoặc bài tập.
    \item Xây dựng báo cáo tiến độ và cảnh báo nguy cơ bỏ học.
\end{itemize}

Khi hệ thống cung cấp tùy chọn tắt cá nhân hóa, người dùng có thể sử dụng tùy chọn này
trong phần cài đặt tài khoản. Khi cá nhân hóa bị tắt, hệ thống vẫn có thể xử lý dữ liệu
tối thiểu cần thiết để cung cấp khóa học, lưu tiến độ và bảo đảm an toàn.

% =========================================================
\section{Quản trị dữ liệu RAG và cơ sở dữ liệu véc-tơ}

Tài liệu được đưa vào hệ thống RAG phải có nguồn gốc và phạm vi sử dụng rõ ràng.

Mỗi nguồn dữ liệu cần được gắn các thông tin quản trị phù hợp, chẳng hạn:

\begin{itemize}
    \item Mã định danh nguồn dữ liệu.
    \item Loại tài liệu và đơn vị sở hữu.
    \item Domain hoặc phân hệ nghiệp vụ.
    \item Vai trò được phép truy cập.
    \item Khóa học, lớp học hoặc tổ chức liên quan.
    \item Phiên bản, thời điểm cập nhật và trạng thái hiệu lực.
\end{itemize}

Trước khi truy xuất dữ liệu RAG, hệ thống phải kiểm tra:

\begin{itemize}
    \item Danh tính và trạng thái tài khoản của người yêu cầu.
    \item Vai trò và quyền hạn hiện tại.
    \item Quyền sở hữu hoặc quyền truy cập khóa học.
    \item Phạm vi domain và các điều kiện bổ sung của nguồn dữ liệu.
\end{itemize}

Dữ liệu véc-tơ chỉ là bản biểu diễn phục vụ truy xuất và không thay thế nguồn dữ liệu gốc.
Khi nguồn gốc bị sửa, thu hồi hoặc hết hiệu lực, dữ liệu liên quan trong Vector DB
phải được cập nhật hoặc xóa theo quy trình đồng bộ.

% =========================================================
\section{Bảo mật và kiểm soát truy cập}

\subsection{Phân quyền RBAC}

Mỗi người dùng chỉ được truy cập các chức năng và dữ liệu cần thiết cho vai trò của mình.

Các vai trò có thể bao gồm:

\begin{itemize}
    \item Học viên.
    \item Giảng viên.
    \item Trợ giảng.
    \item Nhân sự hoặc tư vấn viên.
    \item Quản trị viên.
    \item Quản trị hệ thống.
\end{itemize}

Ngoài việc kiểm tra tại giao diện, mọi API phải kiểm tra quyền ở phía Backend.
Việc ẩn nút hoặc chức năng trên Frontend không được xem là một biện pháp phân quyền đầy đủ.

\subsection{Bảo vệ thông tin xác thực}

\begin{itemize}
    \item Mật khẩu phải được băm bằng thuật toán an toàn như BCrypt.
    \item Mật khẩu không được lưu hoặc ghi log dưới dạng văn bản thuần.
    \item Access Token phải có thời hạn phù hợp.
    \item Refresh Token phải được kiểm soát, thu hồi hoặc vô hiệu hóa khi cần thiết.
    \item Secret key và thông tin kết nối không được ghi trực tiếp trong mã nguồn.
    \item Thông tin nhạy cảm phải được quản lý qua biến môi trường hoặc hệ thống quản lý bí mật.
\end{itemize}

\subsection{Bảo mật dịch vụ nội bộ}

Kết nối giữa Frontend, Backend, AI Service, cơ sở dữ liệu,
Redis, Vector DB và kho tệp phải được giới hạn theo mạng, quyền và môi trường triển khai.

API nội bộ không được công khai nếu không cần thiết. Những API buộc phải công khai
phải có cơ chế xác thực, giới hạn lưu lượng và ghi nhận nhật ký truy cập.

\subsection{Bảo vệ tệp và tài liệu}

\begin{itemize}
    \item Tệp tải lên phải được kiểm tra loại, kích thước và định dạng.
    \item Tên tệp do người dùng cung cấp không được sử dụng trực tiếp làm đường dẫn lưu trữ.
    \item URL truy cập tệp riêng tư phải có thời hạn hoặc yêu cầu xác thực.
    \item Tài liệu khóa học trả phí chỉ được cấp cho người có quyền sở hữu hợp lệ.
\end{itemize}

% =========================================================
\section{Kiểm duyệt nội dung và an toàn AI}

Đầu vào và đầu ra của AI có thể được kiểm tra nhằm phát hiện:

\begin{itemize}
    \item Nội dung kích động bạo lực hoặc tự gây hại.
    \item Nội dung thù ghét, quấy rối hoặc phân biệt đối xử.
    \item Nội dung tình dục không phù hợp.
    \item Nội dung hướng dẫn hành vi nguy hiểm hoặc vi phạm pháp luật.
    \item Dữ liệu cá nhân hoặc thông tin bí mật có nguy cơ bị tiết lộ.
    \item Prompt injection hoặc yêu cầu vượt qua chính sách hệ thống.
\end{itemize}

Tùy theo mức độ rủi ro, hệ thống có thể từ chối yêu cầu, giới hạn câu trả lời,
đưa ra cảnh báo, chuyển yêu cầu cho người phụ trách hoặc khóa tạm thời chức năng liên quan.

% =========================================================
\section{Chất lượng, sai lệch và độ chính xác của AI}

Nội dung do AI tạo ra có thể chứa lỗi, thiếu ngữ cảnh hoặc phản ánh sai lệch của dữ liệu nguồn.

AILMS áp dụng các biện pháp giảm thiểu:

\begin{itemize}
    \item Ưu tiên truy xuất từ nguồn tài liệu được kiểm duyệt.
    \item Hiển thị nguồn tham khảo khi chức năng hỗ trợ.
    \item Giới hạn AI trả lời trong phạm vi dữ liệu và vai trò được cấp quyền.
    \item Cho phép người dùng đánh giá câu trả lời.
    \item Định kỳ đánh giá chất lượng, độ chính xác và mức độ an toàn của hệ thống.
    \item Yêu cầu giảng viên kiểm tra nội dung AI trước khi sử dụng làm tài liệu chính thức.
\end{itemize}

% =========================================================
\section{Audit Log và giám sát hệ thống}

AILMS ghi nhận Audit Log đối với các hoạt động có mức độ ảnh hưởng cao, bao gồm:

\begin{itemize}
    \item Đăng nhập, đăng xuất và thay đổi trạng thái tài khoản.
    \item Thay đổi vai trò hoặc quyền hạn.
    \item Tạo, sửa, xóa khóa học và nội dung quan trọng.
    \item Thay đổi điểm, kết quả hoặc trạng thái học tập.
    \item Truy cập, đồng bộ hoặc xóa nguồn dữ liệu RAG.
    \item Thực hiện hành động quản trị thông qua AI.
    \item Thay đổi cấu hình bảo mật hoặc tích hợp hệ thống.
\end{itemize}

Audit Log phải được bảo vệ khỏi việc sửa đổi trái phép và chỉ được truy cập bởi
người có nhiệm vụ phù hợp. Không ghi mật khẩu, secret key, access token đầy đủ
hoặc dữ liệu nhạy cảm không cần thiết vào nhật ký.

% =========================================================
\section{Thời hạn lưu trữ và xóa dữ liệu}

Dữ liệu chỉ được lưu trong thời gian cần thiết để thực hiện mục đích đã công bố,
đáp ứng yêu cầu vận hành hoặc tuân thủ nghĩa vụ pháp lý.

Thời hạn cụ thể có thể khác nhau tùy nhóm dữ liệu:

\begin{itemize}
    \item Dữ liệu tài khoản được lưu trong thời gian tài khoản hoạt động và thời gian cần thiết sau khi đóng tài khoản.
    \item Dữ liệu học tập được lưu theo yêu cầu đào tạo và xác nhận kết quả.
    \item Lịch sử hội thoại AI được lưu theo cấu hình và mục đích hỗ trợ người dùng.
    \item Nhật ký bảo mật được lưu trong khoảng thời gian cần thiết để điều tra và kiểm toán.
    \item Tệp tạm và dữ liệu đệm phải được tự động xóa khi hết thời hạn sử dụng.
\end{itemize}

Khi hết thời hạn lưu trữ, dữ liệu phải được xóa, ẩn danh hoặc tổng hợp
theo quy trình kỹ thuật phù hợp.

% =========================================================
\section{Quyền của người dùng đối với dữ liệu}

Tùy theo quy định áp dụng và phạm vi hệ thống hỗ trợ, người dùng có thể yêu cầu:

\begin{itemize}
    \item Được thông báo về dữ liệu đang được thu thập và xử lý.
    \item Truy cập thông tin cá nhân và dữ liệu học tập của mình.
    \item Điều chỉnh thông tin không chính xác.
    \item Yêu cầu xóa dữ liệu khi không còn căn cứ lưu trữ hợp lệ.
    \item Tắt hoặc hạn chế một số tính năng cá nhân hóa.
    \item Phản đối hoặc yêu cầu xem xét quyết định có sự hỗ trợ của AI.
    \item Báo cáo câu trả lời AI không chính xác hoặc không phù hợp.
\end{itemize}

Một số yêu cầu có thể bị giới hạn nếu dữ liệu cần được lưu để bảo vệ quyền lợi hợp pháp,
ngăn chặn gian lận, xác nhận kết quả học tập hoặc tuân thủ nghĩa vụ pháp lý.

% =========================================================
\section{Chia sẻ dữ liệu và dịch vụ bên thứ ba}

AILMS chỉ chia sẻ dữ liệu với bên thứ ba trong các trường hợp cần thiết và có căn cứ phù hợp,
chẳng hạn như:

\begin{itemize}
    \item Nhà cung cấp hạ tầng lưu trữ hoặc điện toán.
    \item Nhà cung cấp dịch vụ AI hoặc mô hình ngôn ngữ.
    \item Đơn vị thanh toán hoặc gửi thông báo.
    \item Cơ quan nhà nước có thẩm quyền theo quy định pháp luật.
\end{itemize}

Dữ liệu được chia sẻ phải giới hạn ở mức cần thiết. Các bên xử lý dữ liệu phải đáp ứng
yêu cầu phù hợp về bảo mật, mục đích sử dụng, thời hạn lưu trữ và xử lý sự cố.

Không được gửi mật khẩu, token truy cập hoặc dữ liệu bí mật không cần thiết
tới nhà cung cấp AI bên thứ ba.

% =========================================================
\section{Ứng phó sự cố an toàn thông tin}

Khi phát hiện hoặc nghi ngờ có sự cố, AILMS có thể thực hiện:

\begin{enumerate}
    \item Ghi nhận, phân loại và đánh giá mức độ ảnh hưởng.
    \item Cô lập dịch vụ, tài khoản hoặc thành phần bị ảnh hưởng.
    \item Thu hồi Access Token, Refresh Token và các phiên đăng nhập liên quan.
    \item Tạm ngừng chức năng AI, RAG hoặc tích hợp có nguy cơ.
    \item Bảo toàn Audit Log và bằng chứng kỹ thuật.
    \item Khắc phục lỗ hổng, phục hồi dữ liệu và kiểm tra tính toàn vẹn.
    \item Thông báo cho các bên bị ảnh hưởng hoặc cơ quan có thẩm quyền khi cần thiết.
    \item Phân tích nguyên nhân và triển khai biện pháp phòng ngừa tái diễn.
\end{enumerate}

Việc đóng băng tài khoản hoặc thu hồi token phải dựa trên mức độ rủi ro và
không được thực hiện tùy tiện nếu chưa có căn cứ kỹ thuật hoặc nghiệp vụ phù hợp.

% =========================================================
\section{Trách nhiệm của người dùng}

Người dùng có trách nhiệm:

\begin{itemize}
    \item Bảo vệ thông tin đăng nhập và không chia sẻ tài khoản.
    \item Cung cấp thông tin chính xác trong phạm vi cần thiết.
    \item Không tải lên tài liệu mà mình không có quyền sử dụng.
    \item Không sử dụng AI để gian lận học tập hoặc xâm phạm quyền của người khác.
    \item Kiểm tra nội dung do AI tạo ra trước khi sử dụng.
    \item Thông báo kịp thời khi phát hiện truy cập trái phép hoặc sự cố.
\end{itemize}

Người dùng có thể bị giới hạn hoặc đình chỉ quyền truy cập nếu vi phạm chính sách,
gây nguy hiểm cho hệ thống hoặc ảnh hưởng đến quyền lợi của người dùng khác.

% =========================================================
\section{Trách nhiệm của giảng viên và quản trị viên}

Giảng viên và quản trị viên phải:

\begin{itemize}
    \item Chỉ truy cập dữ liệu phục vụ công việc được giao.
    \item Không sao chép hoặc chia sẻ dữ liệu học viên trái phép.
    \item Kiểm tra nội dung do AI tạo ra trước khi công bố.
    \item Không dùng dữ liệu nội bộ để huấn luyện hoặc thử nghiệm mô hình bên ngoài khi chưa được phép.
    \item Bảo đảm tài liệu đưa vào RAG có nguồn gốc, quyền sử dụng và phạm vi truy cập rõ ràng.
    \item Báo cáo sự cố hoặc hành vi bất thường cho đơn vị phụ trách.
\end{itemize}

% =========================================================
\section{Quản lý thay đổi và đánh giá định kỳ}

Các mô hình AI, prompt hệ thống, công cụ quản trị và nguồn dữ liệu RAG phải được
đánh giá trước khi triển khai vào môi trường chính thức.

Việc đánh giá có thể bao gồm:

\begin{itemize}
    \item Kiểm tra độ chính xác và chất lượng câu trả lời.
    \item Kiểm tra phân quyền và khả năng rò rỉ dữ liệu.
    \item Kiểm tra prompt injection và hành vi vượt giới hạn.
    \item Kiểm tra hiệu năng, chi phí và giới hạn lưu lượng.
    \item Kiểm tra khả năng truy vết và khôi phục khi có lỗi.
\end{itemize}

Mọi thay đổi quan trọng phải được ghi nhận, kiểm thử và có phương án quay lui phù hợp.

% =========================================================
\section{Sửa đổi chính sách}

AILMS có thể cập nhật chính sách này để phản ánh thay đổi về chức năng,
công nghệ, mô hình AI, quy trình vận hành hoặc quy định pháp luật.

Phiên bản mới phải ghi rõ ngày hiệu lực. Đối với thay đổi ảnh hưởng đáng kể
đến quyền lợi hoặc cách xử lý dữ liệu, AILMS sẽ thực hiện biện pháp thông báo phù hợp.

% =========================================================
\section{Thông tin liên hệ}

Người dùng có thể gửi yêu cầu, phản ánh hoặc báo cáo sự cố thông qua:

\begin{itemize}
    \item \textbf{Đơn vị phụ trách:} Phòng Công nghệ và Dữ liệu AILMS.
    \item \textbf{Email hỗ trợ:} \href{mailto:support@ailms.example}{support@ailms.example}.
    \item \textbf{Kênh hỗ trợ:} Chức năng hỗ trợ hoặc gửi yêu cầu trên hệ thống AILMS.
\end{itemize}

\textit{Lưu ý: địa chỉ thư điện tử trên cần được thay thế bằng địa chỉ chính thức
trước khi ban hành tài liệu.}

% =========================================================
\section{Hiệu lực thi hành}

Chính sách này có hiệu lực kể từ ngày được công bố và áp dụng cho toàn bộ
người dùng, nhân sự và hệ thống thuộc phạm vi quản lý của AILMS.

Các đơn vị liên quan có trách nhiệm phổ biến, triển khai và định kỳ rà soát
việc tuân thủ chính sách.

\vfill

\begin{center}
    \textbf{ĐẠI DIỆN PHÒNG CÔNG NGHỆ VÀ DỮ LIỆU AILMS}\\[1.5cm]
    \textit{Ký và ghi rõ họ tên}
\end{center}

\end{document}